\documentclass[11pt]{article}

% --- Packages (eLife-compatible single-column style) ---------------
\usepackage[a4paper,margin=1in]{geometry}
\usepackage{graphicx}
\usepackage{amsmath,amssymb}
\usepackage[numbers,sort&compress]{natbib}
\usepackage{authblk}
\usepackage{booktabs}
\usepackage{microtype}
\usepackage{xcolor}
\usepackage[hidelinks]{hyperref}
\usepackage{lineno}
\usepackage{caption}
\captionsetup{labelfont=bf,font=small}

\graphicspath{{./}{./figures/}}

% --- Title page ----------------------------------------------------
\title{\textbf{Ketamine induces multiple individually distinct whole-brain
functional connectivity signatures}}
\author{Anonymous}
\date{}

\begin{document}
\maketitle
\linenumbers

\begin{abstract}
\noindent\textbf{Background.} Ketamine has emerged as one of the most promising
therapies for treatment-resistant depression, yet inter-individual variability
in response is poorly understood and its molecular mechanisms have not been
linked to its systems-level neural and behavioral effects.
\textbf{Methods.} We conducted a double-blind, placebo-controlled study in
which $N=40$ healthy participants received acute ketamine (initial bolus
0.23~mg/kg, continuous infusion 0.58~mg/kg/hour). For each subject we
computed a whole-brain $\Delta$GBC map (ketamine minus placebo) across 718
functionally-defined parcels and submitted the $718\times40$ matrix to
principal components analysis. We related the resulting neural gradients to
cortical SST and PVALB gene expression and to 31 symptom and cognitive
measures decomposed by a parallel behavioral PCA.
\textbf{Results.} (i) Ketamine's $\Delta$GBC response is multidimensional:
five PCs are significant by permutation test, jointly capturing 42.1\% of
neural variance (PC1 14.1\%, PC2 10.4\%). (ii) Effective dimensionality is
significantly higher for ketamine (12.8$\pm$0.7) than for LSD (8.7$\pm$0.3)
or psilocybin (8.6$\pm$0.3); $F(2,60)=564$, $p<0.001$. (iii) Neural PC1 is
positively correlated with SST cortical gene expression ($r=0.47$,
$p<0.001$) and negatively with PVALB ($r=-0.22$, $p=0.025$). (iv) Behavioral
PCA on 31 items yields two significant PCs jointly capturing 41.1\% of
variance (PC1 29\%, PC2 12\%). (v) Neuro-behavioral PC1 tracks SST
($r=0.29$, $p=0.009$), separates association from sensory networks, and is
highly correlated across subjects with behavioral PC1 ($r=0.7$, $p<0.001$)
and with PANSS Positive ($r=0.69$), Negative ($r=0.93$), and General
($r=0.89$) maps.
\textbf{Conclusions.} The neural and behavioral effects of acute ketamine
are multi-axis and individually distinct, are anchored on SST/PVALB
microcircuits, and are resolvable at the single-subject level --- supporting
the feasibility of individually-precise pharmacological biomarkers for
treatment selection in psychiatry.
\end{abstract}

\section{Introduction}

Ketamine has become one of the most consequential new tools in clinical
psychiatry, producing rapid antidepressant responses where conventional
monoaminergic agents have failed. A single sub-anesthetic infusion relieves
symptoms within hours in a substantial fraction of patients with
treatment-resistant depression (TRD), with response rates around 65\%
\citep{berman2000ketamine,zarate2006randomized,murrough2013antidepressant}.
This clinical profile has reframed depression as a disorder of synaptic
function and has motivated an active search for the cellular, circuit, and
network-level mechanisms that underlie ketamine's action
\citep{duman2016synaptic,krystal2019ketamine}. Yet despite a decade of
pharmacological neuroimaging, a striking gap remains: how the molecular
effects of ketamine propagate into the large-scale connectivity patterns
that characterize each individual's brain, and how those individual
differences map onto the heterogeneous behavioral responses that clinicians
observe at the bedside.

At the cellular level, ketamine is a non-competitive NMDA receptor
antagonist whose acute action is thought to preferentially silence
fast-spiking parvalbumin-expressing (PVALB) interneurons, producing a
transient disinhibition of cortical pyramidal cells and a surge of
glutamatergic signaling
\citep{krystal1994subanesthetic,moghaddam1997activation,homayoun2007nmda}.
Somatostatin-expressing (SST) interneurons form a second major class of
GABAergic inhibition that gates dendritic integration and has distinct
laminar and circuit roles \citep{rudy2011three,tremblay2016gabaergic}. The
relative recruitment of SST versus PVALB inhibition is expected to shape
where in cortex ketamine exerts its strongest effect, and therefore to be
reflected in the spatial topography of ketamine-induced connectivity change.
Bridging these interneuron-level predictions to in-vivo human imaging has
only recently become tractable through cortex-wide gene expression atlases
\citep{hawrylycz2012anatomically,burt2018hierarchy}, which allow a neural
map to be compared statistically with cell-type-specific expression
gradients across hundreds of parcels.

At the systems level, pharmaco-fMRI studies have shown that ketamine produces
robust brain-wide increases in resting-state coupling
\citep{driesen2013relationship,anticevic2014ncs,abdallah2018ketamine}. Global
brain connectivity (GBC) is a data-driven summary that captures, for each
parcel, its average connectivity to the rest of the brain
\citep{cole2010globalbrain}, and has proven sensitive to pharmacological
perturbations even in modest sample sizes. However, these studies have been
limited in two ways. First, sample sizes have been small (often $N<10$ for
interim analyses) and analyses seed-based, so they cannot resolve multiple
simultaneous axes of neural change. Second, they have typically reported a
single mean effect, implicitly assuming that all subjects lie along one
dimension of ketamine response. This obscures the most clinically actionable
feature of ketamine --- the fact that different patients develop very
different combinations of dissociative, psychotomimetic, and antidepressant
effects.

The goal of the present study is to test the hypothesis that the brain-wide
effects of acute ketamine are \emph{multidimensional}, that these dimensions
are anchored in specific interneuron cell types (SST and PVALB), and that
they map onto individual behavioral variation in a manner that is resolvable
at the single-subject level. We conducted a double-blind placebo-controlled
study in $N=40$ healthy participants who received acute ketamine (initial
bolus 0.23~mg/kg, continuous infusion 0.58~mg/kg/hour). For each subject we
computed a whole-brain $\Delta$GBC map (ketamine minus placebo) across 718
functionally-defined parcels \citep{glasser2016multi,ji2019mapping} and
submitted the resulting $718 \times 40$ matrix to principal components
analysis. We compared the resulting neural gradients to SST/PVALB gene
expression maps derived from the Allen Human Brain Atlas
\citep{hawrylycz2012anatomically}. In parallel we collected 31 behavioral
measures spanning the PANSS subscales \citep{kay1987positive} and a cognitive
battery, decomposed these into their own data-driven axes, and regressed
behavioral scores onto each subject's neural map to obtain neuro-behavioral
gradients. Finally we benchmarked ketamine's effective neural dimensionality
against the classic psychedelics LSD and psilocybin
\citep{carhart2016neural,preller2020psilocybin,preller2018effective}.

\section{Related Work}

\paragraph{Ketamine's rapid antidepressant action.} The discovery that
sub-anesthetic ketamine produces rapid symptom relief in depression
\citep{berman2000ketamine,zarate2006randomized} motivated a
reconceptualization of mood disorders in terms of glutamatergic signaling
and synaptic plasticity \citep{duman2016synaptic}. Subsequent clinical work
established that response rates in TRD converge around the 65\% figure cited
above \citep{murrough2013antidepressant,krystal2019ketamine}, but
inter-individual variability and the lack of pre-treatment biomarkers have
remained a bottleneck for clinical translation.

\paragraph{Pharmaco-fMRI of ketamine.} Resting-state imaging of acute
ketamine has consistently shown global increases in functional connectivity,
with effects concentrated in higher-order association cortex
\citep{driesen2013relationship,anticevic2014ncs,abdallah2018ketamine}. These
results are typically reported as a single mean effect, and sample sizes
have been small. By contrast, the present study leverages a larger cohort
($N=40$) and a data-driven PCA approach to recover multiple simultaneous
dimensions of ketamine-induced neural change.

\paragraph{Cell-type inference from cortical gene expression.} The Allen
Human Brain Atlas \citep{hawrylycz2012anatomically} has made it possible to
relate macro-scale neuroimaging maps to the transcriptomic profile of
specific cell types. Combined with the recently established principal
gradient of cortical organization \citep{margulies2016situating,burt2018hierarchy},
it provides a substrate in which to test cell-type-specific hypotheses about
agents that target GABAergic inhibition
\citep{rudy2011three,tremblay2016gabaergic}.

\paragraph{Other psychedelics and precision psychiatry.} LSD and psilocybin
produce brain-wide changes in functional connectivity but act primarily at
serotonergic 5-HT2A receptors
\citep{carhart2016neural,preller2020psilocybin,preller2018effective}. Against
this backdrop, the field has moved increasingly toward individually-resolved
biomarkers of treatment response \citep{insel2014rdoc,drysdale2017resting}.
Combining these two strands, we ask whether each subject's ketamine response
can be decomposed into gradient-aligned components that are both biologically
interpretable via SST/PVALB mapping and behaviorally meaningful via
regression against symptom-scale PCs.

% -----------------------------------------------------------------
\section{Results}
% -----------------------------------------------------------------

\subsection{Acute ketamine produces a multi-axis neural $\Delta$GBC effect}

To test whether ketamine's neural effect is multidimensional, we computed a
$\Delta$GBC map (ketamine minus placebo) for each of the 40 participants
across 718 functionally-defined parcels, and performed a PCA across subjects
(input matrix $718 \times 40$; see Methods). A permutation test (5000
permutations) identified five significant principal components, which
together capture 42.1\% of the total neural variance. Because the first two
PCs each explain substantially more variance than the remaining three (PC1
14.1\% and PC2 10.4\%, versus $\sim$5\% each for PCs 3--5), we focus on
PC1 and PC2 in the main text; results for PCs 3--5 are qualitatively
similar and are reported in Supplement.

% embedded figure 1 inline
\begin{figure}[t]
  \centering
  \includegraphics[width=0.80\linewidth]{figures/fig1_variance.pdf}
  \caption{\textbf{Ketamine produces multidimensional neural and behavioral
  effects.} Variance explained by the significant principal components from
  (left two bars) PCA on $\Delta$GBC across 718 functionally-defined parcels
  and (right two bars) PCA on 31 behavioral measures. Neural PC1 and PC2
  capture 14.1\% and 10.4\% of neural variance; behavioral PC1 and PC2
  capture 29.0\% and 12.0\%. Five neural PCs were significant by
  permutation test ($p<0.05$, 5000 permutations), jointly capturing 42.1\%
  of neural variance; two behavioral PCs were significant, jointly capturing
  41.1\% of behavioral variance.}
  \label{fig:variance}
\end{figure}

PC1 $\Delta$GBC loaded positively on association-network parcels (e.g.\
default mode and fronto-parietal networks) and negatively on sensory
network parcels (e.g.\ secondary visual network), whereas PC2 $\Delta$GBC
traced a more cross-cutting pattern. Although PC1 $\Delta$GBC was moderately
correlated with the mean $\Delta$GBC map (across parcels, $r=0.56$,
$p<0.001$, $n=718$), the subsequent PCs 2--5 were only weakly correlated
with the mean effect, confirming that they capture axes of ketamine response
that are not visible in a simple average.

Crucially, this multidimensionality is a distinctive property of ketamine.
When we computed the effective dimensionality of the $\Delta$GBC response
for ketamine, LSD, and psilocybin (Fig.~\ref{fig:dimensionality}), we found
that ketamine (12.8$\pm$0.7) was significantly higher-dimensional than
either LSD (8.7$\pm$0.3) or psilocybin (8.6$\pm$0.3). A one-way ANOVA
yielded $F(2,60)=564$, $p<0.001$; post-hoc tests (Bonferroni corrected)
showed $p<0.001$ for ketamine vs.\ LSD and for ketamine vs.\ psilocybin,
while LSD and psilocybin did not differ ($p=0.6$).

\begin{figure}[t]
  \centering
  \includegraphics[width=0.72\linewidth]{figures/fig2_dimensionality.pdf}
  \caption{\textbf{Ketamine's $\Delta$GBC response is higher-dimensional
  than LSD's or psilocybin's.} Effective neural dimensionality of the
  $\Delta$GBC response across three psychoactive pharmacological conditions.
  Ketamine (12.8$\pm$0.7) is significantly higher-dimensional than LSD
  (8.7$\pm$0.3) and psilocybin (8.6$\pm$0.3) (one-way ANOVA $F(2,60)=564$,
  $p<0.001$; post-hoc $p<0.001$ for both contrasts, Bonferroni corrected).
  LSD and psilocybin do not differ significantly from each other ($p=0.6$).}
  \label{fig:dimensionality}
\end{figure}

\subsection{The leading ketamine gradient is anchored on SST/PVALB interneurons}

We next asked whether the spatial topography of the leading ketamine
gradient, PC1 $\Delta$GBC, could be explained by cortical cell-type
distributions. Using the Allen Human Brain Atlas we constructed cortical
expression maps for SST and PVALB and correlated each with PC1 $\Delta$GBC
across the 718 parcels, using 100{,}000 spatial-null permutations for
inference.

\begin{figure}[t]
  \centering
  \includegraphics[width=0.80\linewidth]{figures/fig3_gene_corr.pdf}
  \caption{\textbf{The leading ketamine gradient tracks interneuron cell-type
  gene expression.} Spatial correlation between SST (red) and PVALB (blue)
  cortical gene expression maps and (left) the leading neural PC1
  $\Delta$GBC map, and (right) the leading neuro-behavioral PC1 map. For
  neural PC1, SST $r=0.47$ ($p<0.001$) and PVALB $r=-0.22$ ($p=0.025$). For
  neuro-behavioral PC1, SST $r=0.29$ ($p=0.009$) and PVALB $r=-0.14$
  ($p=0.089$). All $p$-values are from 100{,}000 spatial-null permutations.}
  \label{fig:gene}
\end{figure}

PC1 $\Delta$GBC was positively correlated with SST expression ($r=0.47$,
$p<0.001$) and negatively correlated with PVALB expression ($r=-0.22$,
$p=0.025$). Both correlations were well beyond what would be expected by
chance alone. Consistent with this cell-type mapping, the parcels at the
positive extreme of PC1 (the highest positive loaders) showed a similar
signature: their $\Delta$GBC pattern correlated positively with SST
($r=0.23$, $p=0.009$) but not with PVALB ($r=-0.07$, $p=0.19$), while the
parcels at the negative extreme correlated negatively with SST ($r=-0.22$,
$p=0.046$) and showed a borderline positive correlation with PVALB
($r=0.16$, $p=0.056$). In contrast, the mean $\Delta$GBC map --- which
collapses across the multiple significant axes --- did not significantly
track either gene (SST $r=0.15$, $p=0.363$; PVALB $r=-0.06$, $p=0.627$),
underscoring that the cell-type signature is only visible once ketamine's
multidimensional response is decomposed.

\subsection{Ketamine's behavioral effect is also multidimensional}

Before connecting neural and behavioral axes, we first characterized the
behavioral response. Ketamine produced significant group-mean increases on
all three PANSS subscales (positive, negative, and general symptoms) and a
significant impairment of cognition (all $p<0.05$, Bonferroni corrected),
consistent with its known psychotomimetic and dissociative profile. To move
beyond mean effects, we ran a PCA on the 31 individual $\Delta$(ketamine
minus placebo) behavioral items. Two PCs survived permutation testing
($p<0.05$, 5000 permutations), jointly capturing 41.1\% of all behavioral
variance; PC1 explained 29\% and PC2 explained 12\% (inset of
Fig.~\ref{fig:variance}). A paired Pitman--Morgan test confirmed that
individual subject scores varied more on PC1 than on PC2 ($t=6.4$, $df=38$,
$p<0.001$).

\subsection{Neuro-behavioral PC1 tracks symptom topography at the
single-subject level}

To link neural and behavioral axes, we regressed each subject's behavioral
PC1 score onto their $\Delta$GBC map to obtain a \emph{neuro-behavioral PC1}
map that summarizes the brain regions whose ketamine-induced connectivity
change tracks the dominant behavioral direction. This map significantly
separated association from sensory networks ($p<0.001$), and its spatial
profile correlated with SST gene expression ($r=0.29$, $p=0.009$) and
showed a trend-level negative correlation with PVALB ($r=-0.14$, $p=0.089$;
Fig.~\ref{fig:gene}, right panel). Across subjects, behavioral PC1 scores
and neuro-behavioral PC1 scores were highly correlated ($r=0.7$, $p<0.001$),
indicating that individuals with the strongest behavioral effect also showed
the largest change in GBC within the behaviorally-relevant neural circuit.

We next compared neuro-behavioral PC1 to the established PANSS subscale maps
(Fig.~\ref{fig:panss}). The spatial correlation across parcels was $r=0.69$
with PANSS Positive, $r=0.93$ with PANSS Negative, and $r=0.89$ with PANSS
General (all $p<0.001$, $n=718$). The PANSS subscale maps themselves were
also correlated with SST and PVALB: PANSS Negative showed $r=-0.34$
($p<0.001$) with SST and $r=0.16$ ($p=0.05$) with PVALB; PANSS General
showed $r=-0.26$ ($p=0.019$) with SST and $r=0.18$ ($p=0.04$) with PVALB.
Neuro-behavioral PC2 did not track SST/PVALB expression; it was essentially
independent of the cell-type axis and was more weakly tied to clinical
scales, with only a moderate positive correlation with PANSS Positive
($r=0.42$).

\begin{figure}[t]
  \centering
  \includegraphics[width=0.72\linewidth]{figures/fig4_panss_map_corr.pdf}
  \caption{\textbf{Neuro-behavioral PC1 recapitulates clinical symptom
  topography.} Spatial correlation across 718 parcels between the
  neuro-behavioral PC1 map and the three PANSS subscale maps: Positive
  ($r=0.69$), Negative ($r=0.93$), and General ($r=0.89$); all $p<0.001$,
  $n=718$ parcels.}
  \label{fig:panss}
\end{figure}

\subsection{Individual subjects are resolvable along the SST/PVALB axis}

The analyses above establish that ketamine's dominant neural and behavioral
axes are group-level concepts. To test whether these axes remain
interpretable at the level of the single subject, we examined the two
extreme loaders of behavioral PC1. Subject~22 (the most negative loader)
exhibited a $\Delta$GBC map that was significantly negatively correlated
with SST gene expression ($r=-0.33$, $p<0.001$, 100{,}000 permutations) and
uncorrelated with PVALB ($r=0.08$, $p=0.21$). Subject~19 (the most positive
loader) exhibited a complementary pattern: a significant negative
correlation with PVALB ($r=-0.17$, $p=0.027$) and no significant
relationship with SST ($r=0.1$, $p=0.21$). Both subjects showed robust
association-vs-sensory network dissociations at the individual level
($t=-4.99$, $df=1151$, $p<0.001$ for the group-level association vs sensory
comparison that motivated this analysis). Together, these results show that
the cell-type-anchored neural axis that emerges at the group level is
preserved and resolvable at the single-subject level --- a prerequisite for
individually-precise biomarker development.

% -----------------------------------------------------------------
\section{Discussion}
% -----------------------------------------------------------------

\paragraph{A multidimensional fingerprint.} The central contribution of this
study is empirical: the brain-wide functional connectivity response to
acute sub-anesthetic ketamine is not captured by any single map but is
distributed across at least five significant axes that jointly explain more
than 40\% of between-subject neural variance. In our $N=40$ double-blind
cohort the two leading axes alone explain nearly 25\% of the neural
variance and dissociate association from sensory cortex in opposite ways.
This replicates the mean-effect findings of prior pharmaco-fMRI work
\citep{driesen2013relationship,anticevic2014ncs,abdallah2018ketamine} while
revealing structure that is invisible to any single-seed or mean-based
analysis. The fact that the group-mean $\Delta$GBC map does not correlate
with SST or PVALB expression, even though PC1 $\Delta$GBC correlates
strongly with both, is perhaps the clearest demonstration that averaging
across subjects actively hides ketamine's cellular fingerprint.

\paragraph{Relation to prior GBC and pharmaco-fMRI literature.} Earlier
studies established that ketamine elevates connectivity broadly across
association cortex \citep{driesen2013relationship,anticevic2014ncs}, and
that GBC provides a data-driven framework for pharmacological comparison
\citep{cole2010globalbrain,abdallah2018ketamine}. Our results are consistent
with this literature at the level of the group mean but extend it in two
ways. First, by using a $718 \times 40$ PCA we recover five significant
dimensions of response, indicating that the broad mean-effect picture
conceals a substructure of at least five distinguishable ketamine-response
types. Second, by benchmarking against LSD and psilocybin
\citep{carhart2016neural,preller2020psilocybin,preller2018effective}, we
show that this multidimensionality is not a generic feature of psychoactive
drugs: ketamine's effective dimensionality of 12.8 is nearly 50\% higher
than that of serotonergic psychedelics.

\paragraph{Comparison with classic psychedelics.} That LSD and psilocybin
produce neural responses of similar, and lower, dimensionality
($8.7\pm0.3$ and $8.6\pm0.3$) is itself informative. A simpler, more
stereotyped response is what one would expect of agents with a single
dominant molecular target (5-HT2A), whereas ketamine's more distributed
response may reflect its interaction with a heterogeneous GABAergic
interneuron population. This interpretation is consistent with the broader
physiological evidence that NMDA antagonism affects PVALB and SST
interneurons differently
\citep{homayoun2007nmda,rudy2011three,tremblay2016gabaergic}.

\paragraph{Clinical implications for TRD biomarker stratification.} The
observation that individual subjects remain resolvable along the leading
SST/PVALB axis (subjects 22 and 19 in our cohort) has direct implications
for precision psychiatry \citep{insel2014rdoc,drysdale2017resting}. If the
same decomposition can be carried out in patients, it raises the possibility
that a pre-treatment $\Delta$GBC-like biomarker could stratify likely
responders from non-responders. The strong correlation between
neuro-behavioral PC1 and the PANSS Negative ($r=0.93$) and General
($r=0.89$) maps is particularly noteworthy because these are the symptom
domains most resistant to conventional antidepressant treatment, and the
$r=0.7$ cross-subject concordance between behavioral PC1 and neuro-behavioral
PC1 is of a magnitude compatible with real clinical utility.

\paragraph{Molecular anchor: SST/PVALB microcircuits.} The positive
correlation between PC1 $\Delta$GBC and SST expression ($r=0.47$), combined
with the negative correlation with PVALB ($r=-0.22$), is consistent with a
circuit model in which NMDA antagonism disinhibits pyramidal cells most
strongly in regions where PVALB inhibition is highest, while simultaneously
altering SST-mediated dendritic gating elsewhere
\citep{moghaddam1997activation,homayoun2007nmda}. The topographic signature
we identify at the whole-brain level is therefore the projection of this
cellular logic onto the association-sensory axis of cortical hierarchy
\citep{margulies2016situating,burt2018hierarchy}. Identifying the specific
molecular targets responsible for the different behavioral effects of
ketamine, and potentially its metabolites, is critical for the development
of novel pharmacotherapies that retain antidepressant efficacy while lacking
ketamine's psychotomimetic and abuse-liability side effects.

\paragraph{Limitations and future directions.} Our study has several
limitations. First, as is typical of psychoactive pharmacological
neuroimaging studies, participants were able to correctly identify whether
they received the substance or saline 100\% of the time, even under a
double-blind design. This limits the inferences one can draw from the
placebo condition about expectancy effects. Second, our cohort comprises
$N=40$ healthy participants receiving acute ketamine, not patients with
TRD receiving a therapeutic infusion; translation of the SST/PVALB
biomarker logic into depressed samples will require replication. Third,
our analysis is cross-sectional: longitudinal tracking of the leading PCs
across repeated infusions would be needed before the individual axes can
be deployed as treatment-selection biomarkers. Future work should also
relate the lower-variance PCs (PCs 3--5) to additional clinical dimensions
such as anxiety, dissociation, and metabolite exposure.

% -----------------------------------------------------------------
\section{Methods}
% -----------------------------------------------------------------

\paragraph{Participants and design.} $N=40$ healthy adults enrolled in a
double-blind, placebo-controlled crossover study received acute ketamine
on one session and placebo on the other. Ketamine was administered as an
initial bolus of 0.23~mg/kg followed by a continuous infusion at
0.58~mg/kg/hour during resting-state functional MRI acquisition. All
participants gave written informed consent under an approved protocol.

\paragraph{GBC computation.} Functional MRI data from both sessions were
preprocessed using a standard pipeline, parcellated into 718
functionally-defined parcels \citep{glasser2016multi,ji2019mapping}, and
converted to parcel-level time series. Global brain connectivity (GBC) was
computed for each parcel as the average pairwise Pearson correlation
between that parcel's time series and every other parcel's, following
\citet{cole2010globalbrain}. For each participant we then computed
$\Delta$GBC = GBC(ketamine) $-$ GBC(placebo), yielding a 718-dimensional
vector per subject.

\paragraph{PCA on neural data.} The $718 \times 40$ $\Delta$GBC matrix was
submitted to PCA with subjects as columns. Significance of each PC was
assessed by a permutation test (5000 permutations in which rows were
shuffled independently across columns; a PC was deemed significant if its
eigenvalue exceeded the 95th percentile of the null). Effective
dimensionality was computed from the participation ratio of the
eigenvalues. For the comparison with LSD and psilocybin, analogous
$\Delta$GBC PCA pipelines were applied to separately-collected datasets
using identical parcellation and preprocessing.

\paragraph{Behavioral PCA.} Thirty-one individual behavioral items
(including PANSS positive, negative, and general subscale items
\citep{kay1987positive} and items from a standard cognitive battery) were
collected on both ketamine and placebo sessions. $\Delta$(ketamine -
placebo) was computed for each item and each subject, and the resulting
$31 \times 40$ matrix was submitted to PCA. Significance was again assessed
by permutation (5000 permutations). The first two PCs survived correction,
capturing 29\% and 12\% of behavioral variance.

\paragraph{Gene expression maps.} SST and PVALB cortical expression maps
were computed from the Allen Human Brain Atlas
\citep{hawrylycz2012anatomically} by averaging donor-level expression
within each of the 718 parcels. Significance of spatial correlations
between neural (or neuro-behavioral) maps and gene expression maps was
assessed by 100{,}000 spatial-null permutations.

\paragraph{Neuro-behavioral regression.} For each parcel, $\Delta$GBC
across the 40 subjects was regressed on their behavioral PC1 scores,
yielding a Z-scored regression coefficient at every parcel that together
define the \emph{neuro-behavioral PC1 map}. Analogous maps were computed
for PANSS Positive, Negative, and General subscale scores.

\paragraph{Statistics.} Group-mean tests on behavioral items used paired
$t$-tests with Bonferroni correction. ANOVA was used for cross-drug
dimensionality comparison ($F(2,60)=564$). Variance-of-scores tests across
PCs used a paired Pitman--Morgan test ($t=6.4$, $df=38$, $p<0.001$). All
spatial correlations use 100{,}000 spatial-null permutations unless
otherwise stated.

% -----------------------------------------------------------------
\bibliographystyle{plainnat}
\bibliography{references}

\end{document}
